\documentclass[aspectratio=169, 10pt]{beamer}

% --- Theme and Visual Styling ---
\usetheme{Metropolis}
\usecolortheme{owl} % Dark theme matching the GUI aesthetics

\usepackage[utf8]{inputenc}
\usepackage{amsmath, amssymb}
\usepackage{graphicx}
\usepackage{booktabs}

% --- Presentation Metadata ---
\title{Yet Another Epsilon Meteor Echo Simulator}
\subtitle{Continuous Forward-Scatter Doppler Synthesis, Multi-Harmonic Wind Shear, and Interactive GUI}
\author{\textbf{Antonio M.}}
\institute{International Meteor Organization (IMO) / Radio Observation Group}
\date{International Meteor Conference 2026\\Barcelonnette, France --- September 24--27, 2026}

\begin{document}

% ---------------------------------------------------------------------------
% Slide 1: Title Slide
% ---------------------------------------------------------------------------
\begin{frame}[plain]
    \titlepage
\end{frame}

% ---------------------------------------------------------------------------
% Slide 2: Context & Motivation
% ---------------------------------------------------------------------------
\begin{frame}{Context \& Motivation}
    \begin{itemize}
        \item \textbf{Why simulate overdense radio meteor echoes?}
        \begin{itemize}
            \item Generate controlled synthetic datasets to train automated detection algorithms.
            \item Validate mesospheric wind shear and velocity inversion models.
            \item Provide an accessible, interactive teaching tool for amateur and professional observers.
        \end{itemize}
        \vspace{0.3cm}
        \item \textbf{Limitations of Classical Analytic Models:}
        \begin{itemize}
            \item Static specular cylinder assumptions fail to reflect dynamic mesospheric reality.
            \item Miss key signatures: body Doppler shifts, multi-component splitting, and phase beating.
        \end{itemize}
    \end{itemize}
\end{frame}

% ---------------------------------------------------------------------------
% Slide 3: 3D Forward-Scatter Geometry
% ---------------------------------------------------------------------------
\begin{frame}{3D Forward-Scatter Geometry}
    \begin{columns}
        \begin{column}{0.55\textwidth}
            \begin{itemize}
                \item Parameterized relative to specular center at altitude $h_{\mathrm{spec}}$ ($85\text{--}130\text{ km}$).
                \item Ground baseline distance $d_{\mathrm{Tx-Rx}}$ ranges from $20\text{ to }200\text{ km}$.
                \item 3D trajectory defined by Azimuth ($\alpha$) and Elevation ($\theta$).
                \item Local Doppler shifts depend on projection onto the \textbf{bistatic bisector}:
            \end{itemize}
            \begin{equation*}
                \vec{u}_b(z) = \frac{\hat{u}_t(z) + \hat{u}_r(z)}{\|\hat{u}_t(z) + \hat{u}_r(z)\|}
            \end{equation*}
        \end{column}
        \begin{column}{0.42\textwidth}
            \begin{block}{Vector Mechanics}
                \begin{itemize}
                    \item $\vec{r}_t(z) = \vec{P}(z) - \vec{R}_{\mathrm{Tx}}$
                    \item $\vec{r}_r(z) = \vec{P}(z) - \vec{R}_{\mathrm{Rx}}$
                    \item Path length: $R_{\mathrm{total}}(z) = \|\vec{r}_t\| + \|\vec{r}_r\|$
                \end{itemize}
            \end{block}
        \end{column}
    \end{columns}
\end{frame}

% ---------------------------------------------------------------------------
% Slide 4: Multi-Harmonic Wind Shear & Dynamics
% ---------------------------------------------------------------------------
\begin{frame}{Multi-Harmonic Wind Shear \& Dynamics}
    \begin{itemize}
        \item \textbf{Layered Mesospheric Wind Field Profile:}
        \begin{equation*}
            v_{\mathrm{wind}}(z) = V_{\mathrm{max}} \left[ a_1 \sin(\pi z) + a_2 \sin(2\pi z) + a_3 \sin(3\pi z) \right]
        \end{equation*}
        \item \textbf{Trail Deformation Dynamics:}
        \begin{itemize}
            \item $t < 0.12\text{ s}$: Rapid transient entry phase along trajectory line.
            \item $t > 0.40\text{ s}$: Background wind shear dominates, generating gradients $\frac{dx(z,t)}{dz}$.
        \end{itemize}
        \item \textbf{Specular Reflection Weighting:}
        \begin{equation*}
            W_{\mathrm{spec}}(z, t) = \exp\left( -\frac{(dx/dz)^2}{\sigma_a^2(t)} \right)
        \end{equation*}
    \end{itemize}
\end{frame}

% ---------------------------------------------------------------------------
% Slide 5: Signal Synthesis Engine
% ---------------------------------------------------------------------------
\begin{frame}{Signal Synthesis Engine}
    \begin{columns}
        \begin{column}{0.5\textwidth}
            \textbf{Continuous Phase Integration:}
            \begin{equation*}
                f_D(z, t) = \frac{\vec{V}(z, t) \cdot \vec{u}_b(z)}{\lambda}
            \end{equation*}
            \begin{equation*}
                \Phi(z, t_k) = \Phi(z, t_{k-1}) + 2\pi (f_0 + f_D) dt
            \end{equation*}
            \vspace{0.1cm}
            \textbf{Composite Signal Composite:}
            \begin{equation*}
                S(t) = E(t) \cdot \sum_{z} W_{\mathrm{spec}}(z, t) \sin\left(\Phi(z, t)\right)
            \end{equation*}
        \end{column}
        \begin{column}{0.48\textwidth}
            \begin{block}{Temporal Envelopes $E(t)$}
                \begin{itemize}
                    \item Fast entry rise time ($\tau = 5\text{ ms}$).
                    \item Thermal ambipolar diffusion decay ($\tau = 1.0\text{ s}$).
                    \item Calibrated additive thermal Gaussian noise ($\text{SNR}_{\mathrm{dB}}$).
                \end{itemize}
            \end{block}
        \end{column}
    \end{columns}
\end{frame}

% ---------------------------------------------------------------------------
% Slide 6: Interactive GUI Architecture
% ---------------------------------------------------------------------------
\begin{frame}{Interactive GUI Architecture}
    \begin{itemize}
        \item \textbf{10 Real-Time Control Sliders:} Distance, Altitude, Carrier Freq, Skew Angle, Wind Velocity, Harmonic Ratios ($a_1, a_2, a_3$), Azimuth, and Elevation.
        \item \textbf{Key Technical Robustness Fixes Implemented:}
        \begin{enumerate}
            \item \textit{Aspect-Ratio Distortion}: Decoupled toggle checkboxes into physical-aspect frame axes.
            \item \textit{macOS Tkinter Crashing}: Replaced \texttt{root.destroy()} with \texttt{root.update()} inside OS file dialog event callbacks to preserve master GUI loop safety.
            \item \textit{Spectrogram Log-Zero Warnings}: Implemented matrix power spectrum flooring ($\max(P, 10^{-12})$) prior to decibel scaling.
        \end{enumerate}
    \end{itemize}
\end{frame}

% ---------------------------------------------------------------------------
% Slide 7: Results & Comparison with Observations
% ---------------------------------------------------------------------------
\begin{frame}{Results \& Spectrogram Characteristics}
    \begin{itemize}
        \item \textbf{Head Echo Signature:} Captures early high-velocity Doppler offsets upon meteor entry.
        \item \textbf{Doppler "Snores":} Successfully reproduces symmetrical spectral broadening driven by wind shear deformation.
        \item \textbf{Multi-Pip Splitting \& Beating:} Replicates phase interference across multiple active reflection centers, matching observations from networks like BRAMS and GRAVES receivers.
    \end{itemize}
\end{frame}

% ---------------------------------------------------------------------------
% Slide 8: Summary & Conclusions
% ---------------------------------------------------------------------------
\begin{frame}{Summary \& Future Work}
    \begin{columns}
        \begin{column}{0.5\textwidth}
            \begin{block}{Conclusions}
                \begin{itemize}
                    \item Built a physically accurate, lightweight Python forward-scatter simulator.
                    \item Combines 3D bistatic geometry with multi-harmonic mesospheric shear dynamics.
                    \item Real-time controls with high-dpi audio and plot export features.
                \end{itemize}
            \end{block}
        \end{column}
        \begin{column}{0.5\textwidth}
            \begin{block}{Future Extensions}
                \begin{itemize}
                    \item Add Faraday rotation & polarization dynamics.
                    \item Direct Digital Synthesizer (DDS) raw I/Q baseband export.
                    \item Automated spectral curve fitting against observed echoes.
                \end{itemize}
            \end{block}
        \end{column}
    \end{columns}
\end{frame}

% ---------------------------------------------------------------------------
% Slide 9: Q&A / Thank You
% ---------------------------------------------------------------------------
\begin{frame}[standout]
    Thank You!\\
    \vspace{0.5cm}
    \large Questions \& Discussion
\end{frame}

\end{document}